% Week 08 architecture: our shop and buyer at the close of Round 3. Build: xelatex architecture.tex
\documentclass[11pt]{article}
\usepackage[a4paper,landscape,margin=10mm]{geometry}
\usepackage{kotex}
\usepackage{amsmath}
\usepackage{tikz}
\usetikzlibrary{arrows.meta,calc}
\pagestyle{empty}

\definecolor{llm}{HTML}{E8833A}     % LLM
\definecolor{code}{HTML}{3A7BD5}    % code and probability models
\definecolor{mem}{HTML}{2E9E6B}     % context and memory
\definecolor{tool}{HTML}{8E5CC2}    % tools
\definecolor{mkt}{HTML}{6B6B6B}

% one text size in every box (\small), one for arrow labels (\footnotesize)
\tikzset{
  box/.style={draw=#1, very thick, rounded corners=3pt, fill=#1!8, align=left,
              inner sep=6pt, anchor=north west, font=\small},
  hero/.style={box=code, line width=2.2pt, fill=code!14, inner sep=7pt},
  arr/.style={-{Stealth[length=2.8mm]}, thick},
  fast/.style={-{Stealth[length=3.4mm]}, line width=1.8pt, draw=code},
  slow/.style={-{Stealth[length=3.4mm]}, line width=1.8pt, draw=llm},
  lbl/.style={font=\footnotesize, fill=white, inner sep=1.5pt},
  side/.style={font=\bfseries\small, text=mkt},
}
\newcommand{\when}[1]{\hfill{\footnotesize\color{mkt}#1}}
\newcommand{\haiku}{\textcolor{llm}{\textbf{haiku}}}

\begin{document}
\centering
{\large\bfseries LLM은 말하고, 결정은 코드가 한다}\par\vspace{1mm}
{\small 파란 화살표: 시장 원장 사본에서 매번 다시 계산한 사후분포가 가게와 buyer의 결정에 들어간다.\quad
주황 화살표: manager가 한 시간마다 원장을 읽고 정책 숫자를 고친다.}\par\vspace{5mm}

\begin{tikzpicture}
% columns (cm): shop 0.6-10.1 | shared 10.7-15.2 | buyer 16.0-25.5 ; gutters at 0.15 and 25.95

% ---- market
\node[box=mkt, minimum width=24.9cm, align=center] (market) at (0.6,0) {%
  \textbf{Round 3 시장}\qquad supply 창(20분마다 3분)\quad$\cdot$\quad
  흥정(A2A) $\to$ ACP checkout $\to$ AP2 정산\quad$\cdot$\quad 공개 장부(ledger)};

\node[side, anchor=south west] at (0.6,-1.95) {가게(seller)};
\node[side, anchor=south east] at (25.5,-1.95) {buyer};

% ---- row 1: the two heroes and the ledger copy
\node[hero, text width=9.0cm] (nv) at (0.6,-2.0) {%
  \textbf{현금 제약 newsvendor 매입} (restock.py) \when{supply 창마다}\\[3pt]
  다음 supply 창까지 몇 개가 팔릴지(사후분포)를 보고\\
  재고가 모자랄 때의 손실과 남을 때의 보관 손실을\\
  맞바꿔 SKU마다 매입량을 정함\\[3pt]
  \textbf{1원당 기대 이익}이 큰 개부터, 현금이 남는 만큼};

\node[hero, text width=9.0cm] (pick) at (16.0,-2.0) {%
  \textbf{기대 Saved 최대 가게$\times$SKU + 등급별 상한} \when{주문마다}\\[3pt]
  가게마다 응할 확률, 상한 안에서 살 확률,\\
  고장률을 뺀 물건 가치를 곱한 기대 Saved로 순위\\[3pt]
  상위 3곳에 동시에 묻고, 첫 호가로 흥정할 곳을 정함\\
  상한: 그 등급 가게 마진의 \texttt{min\_share} 이상을 깎는 값};

\node[box=mem, text width=4.0cm] (ledger) at (10.7,-2.0) {%
  \textbf{시장 원장 사본}\\[2pt]
  시장 전체 기록을\\ 라운드 처음부터 복사\\ + 우리 기록\\
  (collector.py)};

% ---- shared posteriors
\node[box=code, text width=4.0cm] (post) at (10.7,-5.0) {%
  \textbf{사후분포} (model.py)\\[2pt]
  카테고리 수요\\ (감마-포아송, 사전은\\ \ 시장 전체)\\
  가게별 체결가 비율\\
  등급별 고장률};

% ---- row 2: per-talk code with haiku
\node[box=code, text width=9.0cm] (seller) at (0.6,-7.2) {%
  \textbf{seller.py} \when{매 흥정}\\[2pt]
  호가: 하한과 시장가 중 큰 값, 식품은 Thompson 밴딧\\
  하한 아래는 거절.\ \ \haiku 는 답장 문장만 씀};

\node[box=code, text width=9.0cm] (review) at (16.0,-7.2) {%
  \textbf{market\_buyer.py 흥정 + review.py} \when{매 주문}\\[2pt]
  제안은 상한까지만 오름.\ \ \haiku 는 문장만 씀\\
  별점: 영수증 사실로 잰 Saved를 목표와 비교};

% ---- row 3: hourly managers and their tools
\node[box=llm, text width=9.0cm] (mgr) at (0.6,-9.8) {%
  \textbf{manager: opus} \when{1시간마다}\\[2pt]
  지난 한 시간을 복기하고 재고 상한·하한 등 숫자를\\
  허용 범위 안에서 고침, 교훈 메모를 다시 씀};

\node[box=tool, text width=4.0cm] (tools) at (10.7,-9.8) {%
  \textbf{툴} (tools.py)\\[2pt]
  sql, lookup\_hour,\\ save\_tool\\
  모두 원장을 읽기만 함};

\node[box=llm, text width=9.0cm] (bmgr) at (16.0,-9.8) {%
  \textbf{buyer manager: opus} \when{1시간마다}\\[2pt]
  지난 주문·흥정·별점을 복기하고 \texttt{min\_share} 등을\\
  허용 범위 안에서 고침, 교훈 메모를 다시 씀};

% ---- fast loop
\draw[fast] (5.4,0 |- nv.north) -- node[lbl,right]{매입} (5.4,0 |- market.south);
\draw[fast, <->] (20.8,0 |- pick.north) -- node[lbl,right]{흥정·결제} (20.8,0 |- market.south);
\draw[fast] (market.south -| ledger.north) -- node[lbl,right]{전부 복사} (ledger.north);
\draw[fast] (ledger.south) -- (post.north);
\draw[fast] (post.west) -| node[lbl,pos=0.75,left]{팔릴 확률·체결가} ($(nv.south east)+(-0.9,0)$);
\draw[fast] (post.east) -| node[lbl,pos=0.75,right]{가격·고장률} ($(pick.south west)+(0.9,0)$);

% ---- slow loop
\draw[slow] (ledger.east) -- ++(0.35,0) |- ($(tools.north east)+(-0.6,0.35)$) -- ($(tools.north east)+(-0.6,0)$);
\draw[slow, <->] (mgr.east) -- (tools.west |- mgr.east);
\draw[slow, <->] (bmgr.west) -- (tools.east |- bmgr.west);
\draw[slow] (mgr.west) -- (0.15,0 |- mgr.west) |- node[lbl,pos=0.25,rotate=90,above]{상한·정책 숫자} ([yshift=-5mm]nv.west);
\draw[slow] (bmgr.east) -- (25.95,0 |- bmgr.east) |- node[lbl,pos=0.25,rotate=90,below]{\texttt{min\_share} 등} ([yshift=-5mm]pick.east);

% ---- buyer forum
\node[box=code, text width=9.0cm] (forum) at (16.0,-12.3) {%
  \textbf{buyer\_forum.py} \when{5분마다}\\[2pt]
  우리 구매를 장부 refs와 함께 사실로만 게시\\
  좋아요를 가장 많이 받은 글 형식을 고름};
\draw[slow] (bmgr.south) -- node[lbl,right]{글 지침(playbook)} (bmgr.south |- forum.north);
\draw[arr] ([yshift=-4mm]ledger.east) ++(0.6,0) |- ($(forum.west)+(0,0)$);

% ---- the second meta-harness layer: the owner and Claude Code
\node[box=mkt, text width=9.0cm] (owner) at (0.6,-12.3) {%
  \textbf{사람 + Claude Code} \when{수시}\\[2pt]
  두 manager의 보고를 텔레그램으로 받아 검토하고,\\
  필요하면 코드 자체를 고침(JOURNAL.md에 기록)};
\draw[arr] (mgr.south) -- node[lbl,right]{텔레그램 보고} (mgr.south |- owner.north);

\end{tikzpicture}

\vfill
\begin{minipage}{24.9cm}\footnotesize
\textcolor{llm}{\rule{3mm}{3mm}} LLM \quad
\textcolor{code}{\rule{3mm}{3mm}} 코드·확률 모델 \quad
\textcolor{mem}{\rule{3mm}{3mm}} 컨텍스트·메모리 \quad
\textcolor{tool}{\rule{3mm}{3mm}} 툴\\[2pt]
키트 기본형: 가게는 haiku가 매 수를 정하고(원가 아래만 금지) restock은 재고가 20개 아래면 무작위 SKU를 산다.
buyer는 재고가 있는 가게 중 무작위 한 곳에서 haiku가 매 제안을 정한다. 원장 사본·사후분포·manager·툴은 없다.
\end{minipage}

\end{document}
